\hspace{3mm}

\centerline{\bf \Huge Теория по степеням вхождения}

\begin{flushright}
    \scriptsize{-Ы Ы Е. Алфавит?\\Куджо Джолин}
\end{flushright}


\textbf{\Large Основные свойства степеней вхождения.}

1) Два натуральных числа $a$ и $b$ равны тогда и только тогда, когда $\nu_p(a) = \nu_p(b)$ для любого простого числа $p$ (иначе говоря, два числа равны тогда и только тогда, когда степени вождения в них всех простых множителей одинаковы)

2)\textit{Важное свойство для решения задач из листочка!!!} Число n делится на число $d$ если и только если любой простой множитель входит в $n$ в не меньшей степени, чем в $d$, т.е. $\nu_p(n)>=\nu_p(d)$ для любого простого $p$; иначе, если для какого-то множителя p это неверно, то при делении образуется дробь с неуничтожаемым множителем p в знаменателе;

3) Последнее условие достаточно проверять, разумеется, только для простых множителей входящих в разложение d; для невходящих будет $\nu_p(d)=0$, что в любом случае не больше $\nu_p(n)$; То есть в задачах где дана дробь вида $\dfrac{a}{d}$ логичнее всего рассматривать простые делители знаменателя, чтобы в числители их было не меньше, а не наоборот.

4) $\nu_p(a)+\nu_p(b) = \nu_p(a\cdot b)$. Это свойство легко проверяется. Пусть $\nu_p(a) = m, \nu_p(b) = n$. Тогда $a = p^md_1$, $b = p^nd_2,$ где $d_1, d_2$ взаимно просты с $p$ потому что мы и так вытащили максимальную степень $p$. Тогда $ab = p^{m+n}d_1d_2$ и тут уже все очевидно.  

5) $\nu_p(a+b)\geq min\{\nu_p(a), \nu_p(b)\}$. Причем равенство будет тогда и только тогда, когда $\nu_p(a) \neq \nu_p(b)$

6) $\nu_p(a^n) = n\nu_p(a)$. Очевидное свойство, докажите сами. 

\hspace{3mm}

\textbf{\large Задачи}

\hspace{3mm}

\problem{1} $a, b, c \ -$ натуральные числа, НОД$(a, b, c) = 1$ и $\dfrac{ab}{a-b} = c$. Докажите, что $a-b$ это точный квадрат.

\hspace{3mm}

\textit{Первое решение.}  Пусть $p \ -$ простой делитель числа $ a - b$. Очевидно, что не может быть так, что $a$ делится, на $p$, но $b$ не делится или наоборот(иначе их разность не делилась бы $p$). Пусть они оба не делятся на $p$, но $a-b$ делится. 
Тогда в дроби $\dfrac{ab}{a-b} = c$ числитель не будет делиться на $p$, когда знаменатель на $p$ делится, а значит эта $p$ в знаменателе ни с кем не сократится и дробь будет не целым числом. Остается один вариант $a \ \vdots \ p$ и $b \ \vdots \ p$. Пусть $\nu_p(a) = k$, $\nu_p(b) = l$, причем $k \neq l$ (не забываем, что $p$ это какой-то простой делитель $a-b$). Тогда можно представить наши числа как $a = p^kd_1, \ b = p^ld_2$. Подставим их в нашу дробь и посмотрим, что получится. 

$\dfrac{ab}{a-b} = \dfrac{p^kd_1p^ld_2}{p^kd_1 - p^ld_2} = \dfrac{p^{k+l}}{p^l(p^{k-l}d_1 - d_2)} = \dfrac{p^k}{p^{k-l}d_1 - d_2}$.

Получаем дробь у которой знаменатель не делится на $p$ (так как мы предположили, что $k \neq l$), а числитель делится на $p$ и дробь в итоге является целым числом. Отсюда очевидно, что это целое число делится $p$(так как в числители эта $p-$шка ни с кем не сократилась.) Но у нас по условию НОД$(a, b, c) = 1$, а они все делятся на $p$. Значит $k = l$, то есть степень вхождения простого числа $p$ (которое делитель числа $a-b$) в числах $a$ и $b$ одинаковая. То есть $\nu_p(a) = \nu_p(b) = k$. Тогда $\nu_p(ab) = 2k$. Но мы знаем, что в итоге ни сверху ни снизу у нас $p-$шек в дроби остаться не могло, так как нельзя числу $c$ давать делиться на $p$. Значит и в знаменателе степень вхождения $p$ точно такая же $2k$, то есть все простые числа в $a-b$ входят в чётной степени, а это и означает, что число является квадратом.

\hspace{3mm}

\textit{Второе решение.} Пусть  НОД$(a, b) = d,  a = dx,  b = dy.$  Поскольку  $ab = ac - bc$,  то $dxy = cx - cy.$  Значит, $cy$ делится на $x$. Так как $x$ и $y$ взаимно просты, $c$ делится на $x$. Аналогично $c$ делится на $y$. Следовательно, c делится на $xy$, т.е.  $c = xyk$.  По условию $d$ и $c$ взаимно просты, тем более d и k взаимно просты. Но  $d = k(x - y)$,  значит, $ k = 1, a - b = dx - dy = d^2.$

\hspace{3mm}

\problem{2} Натуральные числа $m$ и $n$ таковы, что $m^2+n^2+m$ кратно $mn$. Докажите, что $m$ является точным квадратом.

\hspace{3mm}

\textit{Решение.} Пусть $d = (n,m), n = n_1d, \ m = m_1d, \ (m_1, n_1) = 1.$ Левая часть делится на m, значит $n^2$ делится на $m$. Значит $d^2*n_1^2$ делится на $dm_1$ сократим на $d$ и получим $dn_1^2$ делится на $m_1$. Так как $(n_1, d) = 1$, то $d$ делится на $m_1$, доможим оба числа на $d$ и получим, что $d^2$ делится на $m_1d = m$. Итого $d^2$ делится на m.
С другой стороны $m^2+n^2+m$ делится на $mn = m_1n_1d^2$, то есть делятся на $d_2$ первые $2$ слагаемых делятся сами, а значит $m$ делится на $d^2.$ Получаем, что $m \ \vdots \ d^2$ и $d^2 \ \vdots \ m$. Эти две делимости дают $m = d^2$, чтд.

\hspace{3mm}

\problem{3} Тройка целых чисел (x,y,z) является решением уравнения:

\centerline{$y^2z + yz^2 = x^3 + x^2z - 2xz^2$} 

Докажите, что z   является кубом целого числа.

\hspace{3mm}

\textit{Решение.} С виду задача кажется очень страшной, но на деле это совсем не так. Сперва стоит обратить внимание, что все слагаемые у нас "третьей степени", то есть в каждом слагаемом участвуют буквы в количестве трёх штук(например, в $y^2z$ две штуки $y$ и одна штука $z$ и так далее.) Такие уравнения называются \textit{однородными.} Когда мы видим однородное уравнение с целыми числами, то нужно сразу же обращать внимание на НОД. Пусть НОД$(x, y, z) = d$. Тогда $x = dx_1, \ y = dy_1, \ z = dz_1$. Подставим эти значения в наше уравнение и посмотрим, что получилось. 

\centerline{$d^3(y_1^2z_1 + y_1z_1^2) = d^3(x_1^3 x_1^2z_1 - 2x_1z_1^2) = y_1^2z_1 + y_1z_1^2 = x_1^3 x_1^2z_1 - 2x_1z_1^2$}

Получаем то же самое уравнение, что и в начале только теперь у нас тройка чисел $(x_1, y_1, z_1)$ которые взаимнопросты. \textbf{Вывод:} когда видим однородное уравнение в целых числах сразу же можем считать, что все числа взаимнопросты(на Олимпиадах это обязательно нужно доказывать.) Хорошо, теперь можем считать, что $x, y, z$ взаимнопросты. Но что делать дальше? Нам нужно доказать какое-то утверждение про $z$, а в нашем уравнении как раз есть много слагаемых, в которых есть $z$. Давайте перенесем их в одну сторону и вынесем $z$. 

\centerline{$z(y^2 + yz - x^2 + 2xz) = x^3$}

Нам как раз нужно доказать, что $z$ является кубом, а справа он и стоит. Но слева помимо $z$ есть еще множитель. Что же с ним сделать? Если хорошо понимать основную теорему арифметики, то если докажем, что $z$ и $y^2 + yz - x^2 + 2xz$ взаимно просты, то $z$ точно будет кубом(постарайтесь понять это сами, если не получится напишите мне, я объясню). Очевидно, что из $y^2 + yz - x^2 + 2xz$ можно убрать слагаемые $2xz$ и $yz$ так как они делятся на $z$ и никакого вклада в НОД не дадут. Значит нужно посчитать НОД$(z, \ y^2 - x^2)$. Пусть у них есть общий простой делитель равный $p>1$. Тогда получаем, что $z \ \vdots \ p$, а значит и $x^3 \ \vdots \ p$, а так как $p$ простое, то и $x \ \vdots \ p$. Но мы знаем, что $y^2 - x^2$ тоже делится на $p$. Отсюда следует, что $y \ \vdots \ p$. Но мы в начале сказали, что $x, y, z$ взаимнопросты. Противоречие. 

\hspace{3mm}

\problem{4} Решите в натуральных числах уравнение

\centerline{$2^x - 2^y = 2016$}

\hspace{3mm}

\textit{Решение.} Очевидно, что $x>y$. Посмотрим, какая степень вхождения двойки слева $2^x - 2^y = 2^y(2^{x-y} - 1)$. Число в скобке очевидно нечётное, поэтому степень вхождения двойки слева равна $y$. Значит и справа она должна быть $y$. Если разложим 2016 на множители, то увидим, что $2016 = 2^5\cdot3^2\cdot7$. То есть $y=5$. А тут уже несложно подставить $y=5$ в уравнение и найти, что $x=11$. Проверяем $2^{11} - 2^5 = 2048 - 32 = 2016$


\hspace{3mm}

\textbf{\Large Вывод и основные приемы}

\hspace{3mm}

Итого, чтобы решать подобные задачи на делимость одного выражения на другое или же в целом уравнения в целых числах нужно хорошо разбираться в следующих темах:

1) \textbf{Основная теорема арифметики.} Нужно понимать, что у каждого числа есть единственное разложение на степени простых чисел, каждый набор степеней уникален.

2) \textbf{Степень вхождения.} Хорошо пользоваться теми 6 свойствами, которые описаны в начале.

3) \textbf{Уметь пользоваться НОДом и НОКом.} Как вы могли видеть в задачах мы часто прибегали к НОДу, а именно выражали числа $a, b$ как $a = da_1, \ b = db_1$, где d это НОД чисел $a$ и $b$, а $a_1, b_1$ взаимно просты с $d$.

4) \textbf{Пользоваться свойствами взаимной простоты.} Тут мы часто пользовались такими свойствами, как "если $x^3$ делится на простое, то и $x$ делится на это же простое" \ или "если $ab$ делится на $c$ и $(a, c) = 1 $, то $b$ делится на $c$". 

Все эти свойства и приемы нужно себе запомнить и закрепить в памяти. А лучший способ их закрепления это решение задач с листочка. Практика и только практика.